\scanpage{33}
\mspage{15}
% archive sheet number: 32
\textcircled{2} \uline{Le groupe} $S\unsure{\widehat{\widehat{\mathscr E}}|\widehat{\widehat{\mathcal E}}}$.

Par le calcul, il s'explicite comme l'ens des systèmes
\diagram{$(1)\qquad$
\begin{tikzcd}[ampersand replacement=\&,row sep=small,column sep=small]
(p,\arrow[d,phantom,"\in"{rotate=-90}]\&\sigma,\arrow[d,phantom,"\in"{rotate=-90}]\&g_0,\arrow[d,phantom,"\in"{rotate=-90}]\&g_1,\arrow[d,phantom,"\in"{rotate=-90}]\&g_\infty)\arrow[d,phantom,"\in"{rotate=-90}]\&=\widetilde u\\
\widehat{\mathbb Z}^{*}\&\unsure{\mathfrak S_3|\mathscr S_3|G_3}\&\widehat\pi\&\widehat\pi\&\widehat\pi\&
\end{tikzcd}}
satisfaisant la condition
\[
\operatorname{int}(g_{\unsure{\infty|0}})(l_{\sigma(\unsure{\infty|0})}^{p})\cdot\operatorname{int}(g_1)(l_{\sigma(1)}^{p})\cdot\operatorname{int}(g_{\unsure{0|\infty}})(l_{\sigma(\unsure{0|\infty})}^{p})=1\tag*{(2)}
\]
et tels que les trois facteurs de ce produit soient un syst. de générateurs de $\widehat\pi$ [avec la relation génératrice (2)]. Le morphisme $u$ défini par ce \deleted{\unsure{morphi}} système est
\[
\left\{\begin{aligned}
u(l_0)&\unsure{\Longrightarrow|\longmapsto|=|\longrightarrow}\operatorname{int}(g_0)l_{\sigma(0)}^p\\
u(l_1)&=\operatorname{int}(g_1)l_{\sigma(1)}^p\\
u(l_\infty)&=\operatorname{int}(g_\infty)l_{\sigma(\infty)}^p
\end{aligned}\right.\tag*{(3)}
\]
\inserted{ce qui explicite l'homom}
\diagram{$(4)\qquad$
\begin{tikzcd}[ampersand replacement=\&]
S\widehat{\widehat{\mathscr E}}\arrow[r,"p"]\&\widehat{\widehat{\mathscr E}}=\operatorname{Aut}_{\mathrm{lac}}(\widehat\pi)
\end{tikzcd}}
L'homomorphisme canonique
\diagram{$(5)\qquad i:\quad$
\begin{tikzcd}[ampersand replacement=\&]
\widehat\pi\times\widehat{\mathbb Z}^{\{0,1,\infty\}}\arrow[r]\&S\widehat{\widehat{\mathscr E}}
\end{tikzcd}}
s'explicite par
\[
i(g;\alpha_0,\alpha_1,\alpha_\infty)=(1,1,g l_0^{\alpha_0},g l_{\unsure{1|0}}^{\alpha_1},g l_\infty^{\alpha_\infty})\qquad\text{donc}\tag*{(6)}
\]
\[
\deleted{\text{donc}}\quad p\,i(g;\alpha_0,\alpha_1,\alpha_\infty)=\operatorname{int}g.\tag*{(7)}
\]
% slanted
\margin{On retrouve que \unsure{$(p)$|$(4)$} est surjectif, et a comme noyau $i(\widehat\pi)$.}
La multiplication dans $\widehat{\widehat{\mathscr E}}$ est donnée par
\[
\underbrace{(p,\sigma,g_0,g_1,g_\infty)}_{\widetilde u}\underbrace{(p',\sigma',g'_0,g'_1,g'_\infty)}_{\widetilde u'}=(pp',\sigma\sigma',u(g'_0)g_{\sigma'(0)},u(g'_1)g_{\sigma'(1)},u(g'_\infty)g_{\sigma'(\infty)}\tag*{(8)}
\]
% The right edge of (8) does not show a secure final tuple parenthesis.
Cette formule n'est pas, du point de vue du calcul
